\chapter{Metodología}

\section{Caracterización del motor}

Para controlar un motor paso a paso, es fundamental conocer sus características eléctricas y mecánicas. En este caso, se utilizó un motor stepper bifásico modelo 17HS3401S. La información detallada del motor se encuentra en la Figura
(\ref{fig:datasheet_motor})

 A continuación, se listan los parámetros obtenidos:

\begin{table}[htbp]
    \centering
    \caption{Características nominales del motor paso a paso. Fuente: Datasheet del fabricante.\ref{tab:especificaciones_motor}}
    \label{tab:especificaciones_motor}
    \begin{tabular}{|l|c|c|c|}
        \hline
        \textbf{Parámetro} & \textbf{Símbolo} & \textbf{Valor} & \textbf{Unidad} \\ \hline
        Voltaje nominal & $V_{nom}$ & 7.3 V & V \\ \hline
        Corriente por fase & $I_{nom}$ & 1  & A \\ \hline
        Resistencia de fase & $R$ & 3.4 (1 $\pm$ 15 $\%$)  & $\Omega$ \\ \hline
        Inductancia de fase & $L$ & 6 (1 $\pm$ 20 $\%$)  & mH \\ \hline
        Torque de mantenimiento & $\tau_{hold}$ & 330  & mN$\cdot$m \\ \hline
        Torque de detención & $\tau_{dtm}$ & 18  & mN$\cdot$m \\ \hline
        Inercia del rotor & $J$ & 68  & g$\cdot$cm$^2$ \\ \hline
        Ángulo de paso & $\theta_{step}$ & 1.8  & grados \\ \hline
        fases & $N_{phases}$ & 2 & - \\ \hline
    \end{tabular}
\end{table}

Estas propiedades corresponden a las dadas por el fabricante, pero existen ciertas propiedades implicitas que se pueden calcular a partir de estas, como la inductancia mutua, la constante de torque, el número de pasos por revolución, entre otras. 

\begin{equation}
\label{eq:propiedades_implicitas}
\left\{
\begin{aligned}
        L_0 &= \frac{L_q + L_d}{2} \\
        L_2 &= \frac{L_q - L_d}{2} \\
        K_e &= \frac{T_{hold}}{\sqrt{2} \cdot I_{nom}} \\
        N_{steps} &= \frac{360}{\theta_{step}} \\
        P &= 2 \cdot N_{phases} \\
\end{aligned}
\right.
\end{equation}


A partir de las ecuaciones descritas en (\ref{eq:propiedades_implicitas}), se obtuvieron los parámetros necesarios para la configuración del observador, los cuales se resumen en la Tabla \ref{tab:parametros_calculados}.

\begin{table}[htbp]
    \centering
    \caption{Parámetros derivados para el modelo dinámico.}
    \label{tab:parametros_calculados}
    \begin{tabular}{|l|c|c|c|}
        \hline
        \textbf{Parámetro} & \textbf{Símbolo} & \textbf{Valor} & \textbf{Unidad} \\ \hline
        Pasos por revolución & $N_{steps}$ & 200 & steps/rev \\ \hline
        Constante de torque/fuerza & $K_e$ & 0.233 & N$\cdot$m/A \\ \hline
        Inductancia promedio & $L_0$ & 6.0 & mH \\ \hline
    \end{tabular}
\end{table}

Adicionalmente, aun que en el datasheet se describe la inductancia promedio, se sabe que esta no es constante y varía con la posición del rotor. Por lo tanto, se realizó una medición experimental de la inductancia en función de la posición del rotor utilizando un tester LCR. Las pruebas se muestran en la Figura (\ref{fig:medicion_inductancia}). En la siguiente tabla se resumen los valores obtenidos para la inductancia en diferentes posiciones del rotor.

\begin{table}[htbp]
    \centering
    \caption{Tabla de medición de inductancia}
    \label{tab:medicion_inductancia}
    \begin{tabular}{|l|c|c|}
        \hline
        \textbf{Paso} & \textbf{La} & \textbf{Lb}  \\ \hline
        0 & 3.205 mH & 5.15 mH \\ \hline
        1 & 3.24 mH & 4.8 mH \\ \hline
        2 & 5.87 mH & 3.19 mH \\ \hline
        3 & 3.195 mH & 5.56 mH \\ \hline
        4 & 6.4 mH & 3.19 mH \\ \hline
        5 & 8.7 mH & 3.18 mH \\ \hline
        6 & 3.18 mH & 7.58 mH \\ \hline
        7 & 8.71 mH & 3.18 mH \\ \hline
        8 & 3.27 mH & 7.5 mH \\ \hline
    \end{tabular}
\end{table}

Dónde La son las mediciones de inductancia de la fase A y Lb son las mediciones de inductancia de la fase B. Se puede observar que la inductancia varía significativamente con la posición del rotor. Usando las ecuaciones de (\ref{eq:propiedades_implicitas}), se pueden calcular los valores de $L_0$ y $L_2$ para cada posición del rotor, lo que permitirá una modelación más precisa del motor paso a paso.

\begin{equation}
\label{eq:calculo_inductancia}
\left\{
\begin{aligned}
L_q &= \max(L_a,L_b) &= 8.71 [mH] \\ 
L_d &= \min(L_a,L_b) &= 3.18 [mH] \\ 
L_0 &= \frac{L_q + L_d}{2} &= 5.945 [mH] \\ 
L_2 &= \frac{L_q - L_d}{2} &= 2.765 [mH] \\
\end{aligned}
\right.
\end{equation}

Se puede comprobar, que el valor de $L_0$ calculado a partir de las mediciones experimentales es consistente con el valor dado en el datasheet, lo que valida la precisión de las mediciones realizadas. Además, la variación de $L_2$ con la posición del rotor proporciona información valiosa para la modelación dinámica del motor paso a paso, lo que permitirá un control más preciso y eficiente del mismo.